\section{Definition and basic properties}

A main tool in our proof of the central limit theorem
will be that of characteristic functions. The basic
idea will be to show that

\[
\mathbf{E}\left[ g\left( \frac{S_n-n\mu}{\sqrt{n}\sigma}\right)\right] \xrightarrow[n\to\infty]{} \mathbf{E} g(N(0,1))
\]

\noindent for a sufficiently large family of functions $g$. It
turns out that the family of functions of the form

\[
g_t(x) = e^{itx}, \qquad (t\in\R),
\]

\noindent is ideally suited for this purpose. (Here and
throughout, $i=\sqrt{-1}$).

\begin{defn}
The characteristic function of a r.v. $X$, denoted
$\varphi_X$, is defined by

\[
\varphi_X(t) = \mathbf{E}\left(e^{i t X}\right) = \mathbf{E}(\cos(t X)) + i \mathbf{E}(\sin(t X)), \qquad (t\in\R).
\]

\end{defn}

\noindent Note that we are taking the expectation of a
\emph{complex-valued random variable} (which is a kind of
two-dimensional random vector, really). However, the
main properties of the expectation operator
(linearity, the triangle inequality etc.) that hold
for real-valued random variables also hold for
complex-valued ones, so this will not pose too much of
a problem.

Here are some simple properties of characteristic
functions. For simplicity we denote
$\varphi = \varphi_X$ where there is no risk of
confusion.

\begin{enumerate}[label=\arabic*.]
\item $\varphi(0) = \mathbf{E} e^{i\cdot 0\cdot X} = 1$.

\item $\varphi(-t) = \mathbf{E} e^{-i t X} = \mathbf{E}
\left(\overline{e^{it X}}\right) =
\overline{\varphi(t)}$
(where $\overline{z}$ denotes the complex conjugate
of a complex number $z$).

\item $|\varphi(t)| \le \mathbf{E}\left|e^{i t X}\right| =
1$
by the triangle inequality.

\item $|\varphi(t)-\varphi(s)| \le \mathbf{E} \left| e^{i t
X} - e^{isX}\right| = \mathbf{E} \left|e^{i s X}
\left(e^{i(t-s)X}-1\right)\right| = \mathbf{E} \left|
e^{i(t-s)X}-1\right|$.
Note also that
$\mathbf{E} \left|e^{i u X}-1\right|\to 0$ as
$u \downarrow 0$ by the bounded convergence theorem.
It follows that $\varphi$ is a uniformly continuous
function on $\R$.

\item $\varphi_{a X}(t) = \mathbf{E} e^{i a t X} =
\varphi_X(a t)$,\
\ $(a\in\R)$.

\item $\varphi_{X+b}(t) = \mathbf{E} e^{i t (X+b)} = e^{i b
t} \varphi_X(t)$,\
\ $(b\in\R)$.

\item $\textbf{Important:}$ If $X,Y$ are independent then

\[
\varphi_{X+Y}(t) = \mathbf{E}\left(e^{it(X+Y)}\right)
= \mathbf{E}\left(e^{itX} e^{itY}\right) = \mathbf{E}\left(e^{itX}\right) \mathbf{E} \left(e^{itY}\right) = \varphi_X(t) \varphi_Y(t).
\]

\noindent Note that this is the main reason why characteristic
functions are such a useful tool for studying the
distribution of a sum of independent random variables.
\end{enumerate}

\noindent $\textbf{A note on terminology.}$ If $X$ has a
density function $f$, then the characteristic
function can be computed as

\[
\varphi_X(t) = \int_{-\infty}^\infty f_X(x) e^{itx}\,dx.
\]

\noindent In all other branches of mathematics, this would be
called the
$\textbf{Fourier transform}$\footnote{Well, more or less -- it is really the inverse
Fourier transform; but it will be the Fourier
transform if we replace $t$ by $-t$, so that is
almost the same thing.}
of $f$.

So the concept of a characteristic function
generalizes the Fourier transform. If $\mu$ is the
distribution measure of $X$, some authors write

\[
\varphi_X(t) = \int_{-\infty}^\infty e^{itx} d\mu(x)
\]

\noindent (which is an example of a
$\textbf{Lebesgue-Stieltjes integral}$) and call this
the $\textbf{Fourier-Stieltjes transform}$ (or just
the Fourier transform) of the measure $\mu$.
